\documentclass[11pt]{article}
\usepackage{ochamath}
\subjectcolor{2F5DA8}
\setsheet{線形代数・電気系院試補充}{LU・QR・擬似逆行列　演習}{https://university-math-crj.pages.dev/linear-algebra/rank-factorizations/}
\begin{document}
\makesheethead{解答}
自作問題。本文の例題とは数値や設定が異なります。条件・途中式・検算を示してください。
\problem[LUと連立式]
$A=\begin{pmatrix}4&1\\2&3\end{pmatrix},b=(6,8)$ をLU分解で解け。
\begin{ans}
$L=\begin{pmatrix}1&0\\1/2&1\end{pmatrix},U=\begin{pmatrix}4&1\\0&5/2\end{pmatrix}$。前進代入は $y=(6,5)$、後退代入は $x_2=2,x_1=1$。$LU=A$、$Ax=(6,8)$ で検算。$\det A=4(5/2)=10$ も一致する。
\end{ans}
\point{逆行列を作らず三角系を解く。}
\problem[QR最小二乗]
$A=\begin{pmatrix}1&0\\1&1\\0&1\end{pmatrix},b=(2,1,0)$ をQRで最小二乗近似せよ。
\begin{ans}
$q_1=(1,1,0)/\sqrt2,q_2=(-1,1,2)/\sqrt6$、$R=\begin{pmatrix}\sqrt2&1/\sqrt2\\0&\sqrt{3/2}\end{pmatrix}$。$Q^{\mathsf T}b=(3/\sqrt2,-1/\sqrt6)$。後退代入で $x=(5/3,-1/3)$。残差 $(-1/3,1/3,-1/3)$、誤差2乗 $1/3$、$A^{\mathsf T}r=0$。
\end{ans}
\point{三角方程式を解いて残差の直交で確認する。}
\newpage
\problem[階数1の擬似逆]
$A=\begin{pmatrix}2&2\\1&1\end{pmatrix},b=(0,1)$。擬似逆と最小ノルム最小二乗解を求めよ。
\begin{ans}
$A=uv^{\mathsf T},u=(2,1),v=(1,1)$ なので $A^+=vu^{\mathsf T}/10=\frac1{10}\begin{pmatrix}2&1\\2&1\end{pmatrix}$。$x_*=A^+b=(1/10,1/10)$、$Ax_*=(2/5,1/5)$。残差 $(2/5,-4/5)$ は $u$ に直交、誤差2乗4/5。全最小二乗解は $x_1+x_2=1/5$。
\end{ans}
\point{零特異値は逆数にせず0のままにする。}
\problem[条件数]
$A=\operatorname{diag}(2,0.02)$ の $\kappa_2(A),\kappa_2(A^{\mathsf T}A)$ を求めよ。
\begin{ans}
特異値は2と0.02なので条件数100。$A^{\mathsf T}A$ の固有値は4と0.0004、条件数10000。正規方程式を形成すると2ノルム条件数が2乗になる。厳密計算で解が存在することと数値誤差に強いことは異なる。
\end{ans}
\point{条件数は最大と最小の非零特異値の比。}
\newpage
\problem[Moore--Penrose条件]
$A=\operatorname{diag}(3,0)$ の $A^+$ を求め、4つの条件を確認せよ。
\begin{ans}
$A^+=\operatorname{diag}(1/3,0)$。$AA^+=A^+A=\operatorname{diag}(1,0)$ はHermiteな射影。$AA^+A=A$、$A^+AA^+=A^+$ が直接成り立つ。逆行列はないが擬似逆は一意に存在し、像の方向だけを逆に戻す。
\end{ans}
\point{擬似逆と通常の逆行列は区別する。}
\problem[列ベクトル化]
$A=\begin{pmatrix}1&2\\0&3\end{pmatrix},B=\operatorname{diag}(2,4),X=\begin{pmatrix}1&3\\2&4\end{pmatrix}$ で vec の公式を確かめよ。
\begin{ans}
$AX=\begin{pmatrix}5&11\\6&12\end{pmatrix}$、$AXB=\begin{pmatrix}10&44\\12&48\end{pmatrix}$ で列順の vec は $(10,12,44,48)$。$\operatorname{vec}X=(1,2,3,4)$ に $B^{\mathsf T}\otimes A=\operatorname{diag}(2A,4A)$ を掛けても同じ。
\end{ans}
\point{vec は列を積む規約で、右係数は単なる転置。}
\end{document}
